% 第 4 章：经典交流求解器族
% 本章文件由 \input 引入，不含导言区。
\section{经典交流求解器族}\label{sec:classic}

本章覆盖四个经典交流求解器：FDPF 快速解耦、线性化直流潮流、分布式松弛与
自适应孤岛。前两者直接以 \texttt{SolverData} 为输入，复用第~\ref{sec:assembly}~章的
母线分类与 $B'/B''$ 电纳阵装配；后两者以 \texttt{HybridPowerSystem} 为输入，内部各自
重建 SolverData 并委托统一 Newton 核心（第~\ref{sec:newton}~章）完成数值求解，自身
只负责参与因子分摊与分岛调度。四者的公共入口在公共 API 族表（第~\ref{sec:overview}~章）
中列出。

\subsection{快速解耦潮流 FDPF（FDXB 变体）}

\compmeta{powerflow::FDPFSolver}{\srcpath{include/hacdcpf/power\_flow/solvers/fdpf\_solver.hpp}}

\subsubsection*{功能说明与入口}
经典 FDXB 方案快速解耦潮流，\textbf{纯 AC}：不读取 DC 网络与换流器，
\fld{vdc} 输出为空。入口签名（include/hacdcpf/power\_flow/fdpf\_solver.hpp；实现
\srcpath{src/power\_flow/fdpf\_solver.cpp:solve}）：
\begin{quote}\footnotesize
\texttt{PowerFlowResult solve(const SolverData\& data, const PowerFlowOptions\& opt) const;}
\end{quote}
\srcpath{include/hacdcpf/power\_flow/fdpf\_solver.hpp} 是 3 行转发头，规范头在
\texttt{solvers/} 子目录。公共 API \srcpath{solve\_power\_flow\_fdpf}
（\srcpath{src/api/hacdcpf.cpp:solve\_power\_flow\_fdpf}）经 SolverDataCache（见第~\ref{sec:overview}~章）
取得装配数据、施加 ZIP 权重后调用本求解器，再统一做支路潮流回算与结果回投影
（:solve\_power\_flow\_fdpf）。

\subsubsection*{数学模型与算法}
母线分类复用 \srcpath{classify\_ac\_buses}
（\srcpath{include/hacdcpf/power\_flow/pf\_utils.hpp:classify\_ac\_buses}）：首个 SLACK 为参考
（\srcpath{first\_slack\_or\_default}），\textbf{额外 SLACK 按 PV 处理}——
与统一 Newton"额外 SLACK 整体剔除"的约定不同（见第~\ref{sec:newton}~章）。
状态初值取母线记录 \fld{vm\_pu}/\fld{va\_deg}（src/power\_flow/fdpf\_solver.cpp:solve）；
PV 母线电压在 Q 步中不更新（Q 步仅作用 PQ 集合，:solve），即 PV 幅值钉在
母线记录初值，\textbf{不读取发电机 \fld{vg} 整定}。

常量阵 $B'/B''$ 由 \srcpath{build\_susceptance\_matrix} 装配
（实现 \srcpath{src/power\_flow/admittance\_builder.cpp:build\_susceptance\_matrix}，
取 $-\mathrm{Im}(Y^{\mathrm{mod}}_{ij})$），构型由 \texttt{BpBuildFlags} 控制
（\srcpath{include/hacdcpf/assembly/ybus\_builder.hpp:BpBuildFlags}，详见第~\ref{sec:assembly}~章）：
$B'$ 置 $R=0$、忽略充电、变比幅值取 1、相移取 0、不含母线并联（\srcpath{make\_bp\_flags}），
并跳过 $|x|<\texttt{min\_x\_pu}$ 的支路；$B''$ 保留 $R$、充电与变比，相移取 0，
加母线并联 $-B_s/S_{\mathrm{base}}$（\srcpath{make\_bpp\_flags}）。
两阵分别降阶至非 slack 全集与 PQ 集合并用 Eigen SparseLU \textbf{一次性分解}
（src/power\_flow/fdpf\_solver.cpp:solve），分解失败直接返回 \fld{converged}=false（:solve）。

每次迭代（:solve）先以\textbf{稀疏复数形式}计算注入
（\srcpath{compute\_power\_injections}）：$V_i=\mathrm{polar}(V_{m,i},\theta_i)$，
$\mathbf I=Y_{bus}\mathbf V$（稀疏矩阵--向量乘），$S_i=V_i\,\overline{I_i}$，
\textbf{不再生成稠密} $Y_{bus}$：
\begin{align*}
P_i^{\mathrm{calc}}&=\sum_{j}V_iV_j\big(G_{ij}\cos\theta_{ij}+B_{ij}\sin\theta_{ij}\big)\\
Q_i^{\mathrm{calc}}&=\sum_{j}V_iV_j\big(G_{ij}\sin\theta_{ij}-B_{ij}\cos\theta_{ij}\big)
\end{align*}
（与第~\ref{sec:newton}~章失配方程同式）。指定注入由
\srcpath{compute\_specified\_injections}（:compute\_specified\_injections）计算，含 ZIP 电压项
（权重顺序恒功率/恒电流/恒阻抗，见第~\ref{sec:assembly}~章）：\fld{has\_component\_loads}
为真时改用\textbf{逐母线} ZIP 系数（\fld{bus\_zip\_pp}/\fld{\_ip}/\fld{\_zp} 与
\fld{bus\_zip\_pq}/\fld{\_iq}/\fld{\_zq}，配 \fld{pd\_pu}/\fld{qd\_pu}，:compute\_specified\_injections；
由 \srcpath{aggregate\_load\_demand} 装配，\srcpath{src/power\_flow/solver\_data.cpp:aggregate\_load\_demand}），
否则沿用母线级 \fld{pd\_mw}/\fld{qd\_mvar} 加全局 \fld{zip\_pw}/\fld{zip\_qw}
（src/power\_flow/fdpf\_solver.cpp:compute\_specified\_injections）。如实说明：旧版 FDPF 只用全局 ZIP 权重，逐母线口径是本次的行为修正，
使 FDPF 与统一 Newton 的负荷模型对齐。
失配按 FDPF 惯例除以幅值：
\[
\Delta P_i/V_i=\big(P_i^{\mathrm{spec}}-P_i^{\mathrm{calc}}\big)/V_i,\qquad
\Delta Q_i/V_i=\big(Q_i^{\mathrm{spec}}-Q_i^{\mathrm{calc}}\big)/V_i
\]
（src/power\_flow/fdpf\_solver.cpp:solve）。收敛判据为
$\max\big(\lVert\Delta P/V\rVert_\infty,\lVert\Delta Q/V\rVert_\infty\big)<\texttt{tol}$
（:solve）。校正分两步单序执行（每迭代 P 步、Q 步各一次）：
\[
B'\,\Delta\theta=\Delta P/V\;(:solve),\qquad
B''\,\Delta V=\Delta Q/V\;(:solve)
\]
如实说明三处简化与缺陷：
\begin{itemize}
  \item \textbf{无全局化与限值逻辑}：无步长截断、无回退链（对照第~\ref{sec:newton}~章），
        无 PV$\leftrightarrow$PQ 无功限值切换，发电机 Q 限值完全不参与；
  \item \textbf{注入计算已稀疏化}（:compute\_power\_injections）：复杂度 $O(\mathrm{nnz}(Y_{bus}))$，
        旧版的 $Y_{bus}$ 转稠密 $+O(n^2)$ 双重循环已移除（主文档诚实口径相应条目
        随之过时，以本节为准）；
  \item 未收敛时复用同一稀疏注入内核与指定注入函数复算最终残差写入 \fld{residual}
        （跳过 slack 的 P 行与非 PQ 母线的 Q 行，:solve）。
\end{itemize}

\subsubsection*{求解参数}
\begin{paramtable}
\fld{tol} & $\varepsilon$ & pu & $10^{-6}\sim10^{-10}$ & $10^{-8}$ &
  收敛容差（$\Delta P/V$、$\Delta Q/V$ 的 $\ell_\infty$ 范数），
  \texttt{Defaults::kPFTol}（\srcpath{include/hacdcpf/model/defaults.hpp:Defaults}）\\
\fld{fdpf\_max\_iter} & $k_{\max}$ & 次 & 50--5000 & 1000 &
  最大迭代数（\srcpath{include/hacdcpf/power\_flow/power\_flow\_options.hpp:PowerFlowOptions}；
  读于 src/power\_flow/fdpf\_solver.cpp:solve）\\
\fld{robust\_nonlinear.min\_branch\_x\_pu} & $x_{\min}$ & pu & --- & $10^{-20}$ &
  $B'$ 装配的小电抗支路剔除阈值（真实头
  \srcpath{include/hacdcpf/power\_flow/globalization/robust\_nonlinear\_options.hpp:RobustNonlinearOptions}；
  传参于 src/power\_flow/fdpf\_solver.cpp:solve；注意它\textbf{覆盖} \srcpath{make\_bp\_flags}
  自带的 $10^{-6}$ 缺省）\\
\end{paramtable}
全局 ZIP 权重 \fld{zip\_pw}/\fld{zip\_qw} 与 \fld{loss\_model} 由 API 层在装配侧读取
（\srcpath{apply\_zip\_weights}，src/api/hacdcpf.cpp:get\_cached\_solver\_data，调用于 :solve\_power\_flow\_fdpf）；
求解器核在 \fld{has\_component\_loads} 为真时改读逐母线 ZIP 系数
（见上文算法小节），否则回退全局权重。

\subsubsection*{验证与校核}
\begin{itemize}
  \item \textbf{冒烟}：\srcpath{tests/test\_hacdcpf.cpp:TEST\_CASE} 两节点算例断言收敛；
  \item \textbf{稀疏注入与逐母线负荷回归}：\srcpath{tests/test\_helm\_solver.cpp}
        （注册于 \srcpath{tests/CMakeLists.txt}）——「FDPF sparse injection path」
        两节点算例与 Newton 逐母线 $|V|/\theta$ 对照，容差 $2\times10^{-4}$（:make\_two\_bus\_hybrid\_dc\_case）；
        「FDPF uses aggregated component loads」挂显式 \texttt{Load} 组件触发
        \fld{has\_component\_loads} 逐母线 ZIP 路径，与 Newton 对照容差
        $3\times10^{-4}$（:make\_two\_bus\_hybrid\_dc\_case）；
  \item \textbf{与 Newton 对照}：\srcpath{tests/test\_matpower\_cases.cpp:TEST\_CASE}
        case30 下两法均收敛时逐母线 $|V|$ 对照，容差 \texttt{WithinAbs}$=10^{-3}$；
        任一不收敛则 \texttt{SKIP}——\textbf{FDPF 无 Tier2/3 批量回归}，覆盖薄；
  \item 能力合同工具 \srcpath{tools/validate\_case\_capabilities.py} 做
        NR/FDPF/DC 三法互验（见第~\ref{sec:validation}~章）。
\end{itemize}

\subsection{线性化直流潮流（$B\theta=P$）}

\compmeta{powerflow::solve\_ac\_linearized\_dc}{\srcpath{include/hacdcpf/power\_flow/ac\_linearized\_pf.hpp}}

\subsubsection*{功能说明与入口}
经典 DC 潮流：$|V|\equiv1$、忽略电阻与无功，角度方程一次线性求解。入口为自由函数
（include/hacdcpf/power\_flow/ac\_linearized\_pf.hpp:solve\_ac\_linearized\_dc；实现 \srcpath{src/power\_flow/ac\_linearized\_pf.cpp:solve\_ac\_linearized\_dc}）：
\begin{quote}\footnotesize
\texttt{ACLinearizedDCResult solve\_ac\_linearized\_dc(const SolverData\& data);}
\end{quote}
结果结构 \texttt{ACLinearizedDCResult\{va（弧度）, pf\_mw, success\}}
（include/hacdcpf/power\_flow/ac\_linearized\_pf.hpp:ACLinearizedDCResult）。公共 API \srcpath{solve\_ac\_dc\_power\_flow}
（\srcpath{src/api/hacdcpf.cpp:solve\_ac\_dc\_power\_flow}）负责回投影相角与支路功率到原始编号空间
（:solve\_ac\_dc\_power\_flow）。如实说明：GUI 将该入口标为 \texttt{hybrid\_linearized}，但实现
\textbf{只解 AC 相角}，不求解 DC 电压、无功与网损。

\subsubsection*{数学模型与算法}
slack 取首个 SLACK 母线，缺省回退 bus 0（src/power\_flow/ac\_linearized\_pf.cpp:solve\_ac\_linearized\_dc）。对每条在运且
$|x_{pu}|\ge10^{-20}$ 的支路（:solve\_ac\_linearized\_dc）取
\[
b_k=\frac{1}{x_{pu}\cdot t_k}\qquad(t_k=0\text{ 按 }1.0\text{ 处理})
\]
（:solve\_ac\_linearized\_dc），装配 Laplacian $B_{ii}\mathrel{+}=b_k$、$B_{ij}=-b_k$
（:solve\_ac\_linearized\_dc）。移相变压器以注入等效：from 端 $+b_k(-\varphi_k)$、to 端
$-b_k(-\varphi_k)$（:solve\_ac\_linearized\_dc，$\varphi_k$ 为 \fld{shift\_deg} 化弧度）。代码先定义
$P^{\mathrm{shift}}$，再严格按注释“$B\theta+P^{\mathrm{shift}}=P^{\mathrm{spec}}$”
装配右端：
\[
P_i=P_{g,i}-P_{d,i}^{\mathrm{pu}}-P_i^{\mathrm{shift}}
\]
（:solve\_ac\_linearized\_dc；当 \fld{pd\_pu} 维数不匹配时才回退为
$P_{d,i}^{\mathrm{pu}}=\fld{pd\_mw}/\fld{base\_mva}$，恒功率口径，
\textbf{不含 ZIP 电压项}）。剔除 slack 行/列得降阶方程（:solve\_ac\_linearized\_dc），右端计入
slack 相角：
\[
B_{\mathrm{red}}\,\theta_{\mathrm{red}}
=\mathbf P_{\mathrm{red}}-\mathbf B_{\mathrm{red,slack}}\,\theta_{\mathrm{slack}}
\]
（:solve\_ac\_linearized\_dc），Eigen SparseLU 一次求解（:solve\_ac\_linearized\_dc）；分解或回代失败时
\fld{success}=false 静默返回。解出全状态后，代码只在非参考节点计算
$\max_i|(B\theta)_i-P_i|$ 并写入 \fld{residual\_pu}（:solve\_ac\_linearized\_dc），不检查参考节点
平衡。支路有功回算
\[
P_{f,k}=b_k\big(\theta_i-\theta_j-\varphi_k\big)\cdot S_{\mathrm{base}}
\]
（:solve\_ac\_linearized\_dc，MW）。近似边界如实声明：电压幅值、无功、网损均不输出；
$r_{pu}$、充电、并联全部忽略；多 slack 时仅首个被剔除，其余按普通母线计入
$B$（与 Newton/FDPF 的分类约定均不同）。

\subsubsection*{求解参数}
本求解器\textbf{不读取任何} \texttt{PowerFlowOptions} 字段（函数签名无 \texttt{opt}，
include/hacdcpf/power\_flow/ac\_linearized\_pf.hpp:solve\_ac\_linearized\_dc），全部阈值为硬编码常数：
\begin{paramtable}
（内部常数）$|x_{pu}|$ 下限 & --- & pu & --- & $10^{-20}$ &
  低于此值的支路视为零阻抗跳过（src/power\_flow/ac\_linearized\_pf.cpp:solve\_ac\_linearized\_dc）\\
（内部常数）$|t|$ 下限 & --- & pu & --- & $10^{-12}$ &
  低于此值变比按 1.0 处理（src/power\_flow/ac\_linearized\_pf.cpp:solve\_ac\_linearized\_dc）\\
\end{paramtable}

\subsubsection*{验证与校核}
\begin{itemize}
  \item \textbf{冒烟}：\srcpath{tests/test\_hacdcpf.cpp:TEST\_CASE} 两节点算例断言
        相角向量非空；
  \item \textbf{投影空间回归}：\srcpath{tests/test\_sppt\_metamorphic.cpp:TEST\_CASE}
        断言 \fld{va} 尺寸=原始母线数、\fld{pf\_mw} 尺寸=原始支路数、被合并支路
        功率保持 0；
  \item 覆盖薄弱如实说明：无对 MATPOWER \texttt{rundcpf} 的批量数值对照，
        本求解器的正确性主要靠 $B\theta=P$ 公式简单性与上述结构断言背书。
\end{itemize}

\subsection{分布式松弛}

\compmeta{powerflow::DistributedSlackSolver}{\srcpath{include/hacdcpf/power\_flow/solvers/distributed\_slack\_solver.hpp}}

\subsubsection*{功能说明与入口}
把 slack 吸收的不平衡功率按参与因子分摊到多台电源。两个入口
（\srcpath{include/hacdcpf/power\_flow/distributed\_slack\_solver.hpp}；实现
\srcpath{src/power\_flow/distributed\_slack\_solver.cpp:solve\_simplified}）：
\begin{quote}\footnotesize
\texttt{DistributedSlackResult solve\_simplified(const HybridPowerSystem\&, const DistributedSlack\&,\\
\hspace*{2em}const PowerFlowOptions\&) const;}\\
\texttt{DistributedSlackResult solve\_full\_jacobian(...同名参数...) const;}
\end{quote}
配套工厂 \srcpath{create\_participation\_factors}（:create\_participation\_factors）
按 \texttt{capacity}/\texttt{droop}/\texttt{equal} 三种方式生成配置结构
\texttt{DistributedSlack}（四字段定义见
\srcpath{include/hacdcpf/power\_flow/power\_flow\_options.hpp:DistributedSlack}，语义列于本节末）。
公共 API \srcpath{solve\_power\_flow\_distributed\_slack}(\srcpath{\_full})
（\srcpath{src/api/hacdcpf.cpp:solve\_power\_flow\_distributed\_slack}）。

\subsubsection*{数学模型与算法}
\textbf{参与因子}（src/power\_flow/distributed\_slack\_solver.cpp:create\_participation\_factors）：默认参与母线为有实际发电（$P_g>0$）的 SLACK/PV
在运母线（:create\_participation\_factors）。发电聚合与参与列表均使用 AC 母线稳定
\fld{index}，不把 ID 当作 vector 位置；显式列表过滤不在运、不存在和重复 ID。
空集抛 \texttt{invalid\_argument}（:create\_participation\_factors）；原始权重
\[
\alpha_i^{0}=\max(P_{g,i},10^{-4})\;(\text{capacity}),\qquad
\alpha_i^{0}=1/R_i\;(\text{droop},\ R_i>0\text{ 否则取容量}),\qquad
\alpha_i^{0}=1\;(\text{equal})
\]
（:create\_participation\_factors），统一归一化 $\sum_i\alpha_i=1$（:normalize\_factors, :create\_participation\_factors）；单机分摊上限
与该母线机组群的调节裕度挂钩（三种模式同一口径），
$\overline P_i=\max\big(|P_{\min,i}^{\Sigma}-P_{g,i}^{\mathrm{sched}}|,\,
|P_{\max,i}^{\Sigma}-P_{g,i}^{\mathrm{sched}}|\big)/S_{\mathrm{base}}$
（:create\_participation\_factors），写入 \fld{max\_participation\_p}（:create\_participation\_factors）。\srcpath{sanitize\_slack\_cfg}
（:sanitize\_slack\_cfg）：参与列表为空时自动按 capacity 生成；因子数不匹配时全置 1 再归一；
\fld{reference\_bus}=0 时取首个在运 SLACK 母线（无则抛
\texttt{invalid\_argument}，:sanitize\_slack\_cfg），并校验其确为在运 AC 母线（:sanitize\_slack\_cfg）。

\textbf{简化变体} \srcpath{solve\_simplified}（:solve\_simplified）：先做换流器协调预校核
（可阻断，:solve\_simplified），随后 \srcpath{make\_solver\_data} $+$ 统一 Newton 求基准解
（:solve\_simplified，电压经 \srcpath{unproject\_bus\_vector} 回投影，:solve\_simplified）；
基准解收敛后计算\textbf{全部 AC 母线}上的总不平衡
\[
\Delta P_{\Sigma}=\sum_{i\in\text{全部 AC 母线}}\big(P_i^{\mathrm{calc}}-P_i^{\mathrm{spec}}\big)
\]
（\srcpath{compute\_total\_participating\_mismatch\_pu}；遍历 Ybus 非零元，按
$P_i=\sum_{j:Y_{ij}\ne0}V_iV_j[G_{ij}\cos(\theta_i-\theta_j)+B_{ij}\sin(\theta_i-\theta_j)]$
复算，时间 $O(n+\mathrm{nnz}(Y))$、附加向量空间 $O(n)$，不构造稠密 Ybus；
列主序遍历保留各行按 $j$ 递增的累加顺序；
含换流器 AC 注入 \srcpath{converter\_ac\_injection}，:compute\_total\_participating\_mismatch\_pu）。求和取全部已解 AC
母线而非仅参与母线：非 slack 方程在数值容差内相消，剩余即全体 slack 母线提供的
有功；只取参与母线会在物理 slack 不在参与集合时静默返回零（函数内注释，
src/power\_flow/distributed\_slack\_solver.cpp:compute\_total\_participating\_mismatch\_pu）。按
$\Delta P_i=\alpha_i\,\Delta P_{\Sigma}$ 分摊输出（:solve\_simplified）。

AUD-092/096 的注册回归为 \srcpath{tests/test\_review\_regressions.cpp} 与
\srcpath{tests/test\_advanced\_pf.cpp}；可复现局部性能对照为
\srcpath{tools/review\_performance\_benchmark.py}，实测范围与误差见
\srcpath{docs/testing/module\_code\_audit.md}，不把 mismatch 耗时解释为整次 PF 耗时。

\textbf{限值变体} \srcpath{solve\_full\_jacobian}：每轮完整 Newton 求解后，用
\srcpath{distribute\_with\_limits} 迭代裁剪重分摊——越限母线钉在发电机
$P_{\min}/P_{\max}$ 与 \fld{max\_participation\_p} 的交集边界并计入
\fld{hit\_limits}，再把修正写回发电机并重解；剩余失配小于
$\max(10^{-10},\fld{tol})$ pu 时结束，最多 20 次完整求解。

诚实口径（重要）：
\begin{itemize}
  \item 简化变体只做一次 Newton，\fld{distributed\_slack\_p} 是事后分摊；限值变体
        会重调度并反复重解。兼容名称 \texttt{full\_jacobian} 不表示把参与因子作为
        单个增广 Jacobian 未知量；
  \item 失配复算复用统一 AC 注入装配，覆盖组件负荷、ZIP 电压项、VSC/LCC 与
        energy-router 注入；网络功率仍由 Ybus 致密复算，复杂度为 $O(n^2)$；
  \item \fld{reference\_bus} 会实际改写角度参考；输出
        \fld{distributed\_slack\_p} 的单位是系统 \fld{base\_mva} 基准上的 pu，
        不是 MW。
\end{itemize}

\subsubsection*{求解参数}
\begin{paramtable}
\fld{tol} & $\varepsilon$ & pu & $10^{-6}\sim10^{-10}$ & $10^{-8}$ &
  基准 Newton 收敛容差（转发，见第~\ref{sec:newton}~章）\\
\fld{max\_iter} & $k_{\max}$ & 次 & 20--300 & 50 &
  基准 Newton 迭代上限（转发；\texttt{Defaults::kPFMaxIter}，
  \srcpath{include/hacdcpf/model/defaults.hpp:Defaults}）\\
\fld{enable\_converter\_coordination\_check} & --- & 布尔 & --- & false &
  换流器协调预校核开关，阻断级问题直接返回失败
  （src/power\_flow/distributed\_slack\_solver.cpp:solve\_simplified）\\
\fld{loss\_model} & --- & 枚举 & --- & Linear &
  换流器损耗模型（\srcpath{make\_solver\_data} 装配与注入复算共用，:compute\_total\_participating\_mismatch\_pu, :solve\_simplified）\\
\end{paramtable}
\texttt{DistributedSlack} 配置字段（\fld{participating\_buses}、
\fld{participation\_factors}、\fld{reference\_bus}、\fld{max\_participation\_p}）
不属于 \texttt{PowerFlowOptions}，语义见上文算法小节。

\subsubsection*{验证与校核}
\srcpath{tests/test\_advanced\_pf.cpp}（标签 \texttt{[distributed\_slack]}）：
\begin{itemize}
  \item \textbf{因子构造}：capacity/equal 因子归一化误差 $\le10^{-12}$（:TEST\_CASE），
        droop 因子随下垂系数严格递减（:TEST\_CASE）；
  \item \textbf{求解}：3 节点两变体收敛且电压在 $(0.85,1.15)$（:TEST\_CASE）；
        限值 $10^{-12}$ 强制触发 \fld{hit\_limits}（:TEST\_CASE）；
  \item \textbf{MATPOWER 回归}：case9 capacity、case14 两变体一致收敛、case30 droop
        收敛（:TEST\_CASE）；集成用例 case14 下 slack 电压保持 $(0.95,1.10)$
        （:TEST\_CASE）。
\end{itemize}

\subsection{自适应孤岛}

\compmeta{powerflow::AdaptiveSolver}{\srcpath{include/hacdcpf/power\_flow/solvers/adaptive\_solver.hpp}}

\subsubsection*{功能说明与入口}
面向拓扑可能分裂的系统：先分岛，活岛各自独立求解，死岛电压置零。入口签名
（include/hacdcpf/power\_flow/adaptive\_solver.hpp；实现 \srcpath{src/power\_flow/adaptive\_solver.cpp:solve}）：
\begin{quote}\footnotesize
\texttt{AdaptiveSolveResult solve(const HybridPowerSystem\& sys, const PowerFlowOptions\& opt,\\
\hspace*{2em}const std::unordered\_map<int, ReactiveLimit>\& q\_limits = \{\}) const;}
\end{quote}
如实说明：\texttt{q\_limits} 形参已接线——\srcpath{apply\_reactive\_limit\_overrides}
（src/power\_flow/adaptive\_solver.cpp:apply\_reactive\_limit\_overrides）先校验
\fld{base\_mva} 为正、限值有限且 \fld{qmin}$\le$\fld{qmax}，再把限值按
\fld{base\_mva} 折算改写该母线首台在运发电机的 \fld{qmin\_mvar}/\fld{qmax\_mvar}
（同母线其余机组清零）；限值非法或引用的 AC 母线无在运发电机时抛
\texttt{std::runtime\_error}。岛描述结构 \texttt{IslandInfo}
（\srcpath{include/hacdcpf/model/ac\_components.hpp:IslandInfo}）。公共 API
\srcpath{solve\_power\_flow\_adaptive}（\srcpath{src/api/hacdcpf.cpp:solve\_power\_flow\_adaptive}）在收敛后
于全系统 SolverData 上补算支路潮流（:solve\_power\_flow\_adaptive）；
\srcpath{solve\_power\_flow\_islanded}（:solve\_power\_flow\_islanded）为薄包装。

\subsubsection*{数学模型与算法}
\textbf{分岛} \srcpath{detect\_islands}
（\srcpath{src/power\_flow/island\_detector.cpp:detect\_islands}）：AC/DC 母线合并为一张图，
停运与 \texttt{ISOLATED}/\texttt{DC\_ISOLATED} 母线剔除（:detect\_islands）；真实母线编号
映射到数组位置，编号\textbf{不要求连续}（:ac\_node（lambda））。边集涵盖 AC 支路、闭合开关与
断路器、2W/3W 变压器、DC 支路、闭合 DC 断路器、VSC、DCDC 与能量路由器端口链
（:detect\_islands），DFS 求连通分量（:detect\_islands）。活岛判据 \fld{has\_generators} 只看
\textbf{真实在运注入}（发电机 $P_g>10^{-9}$ 或 \fld{is\_slack}、换流器
$|P_{set}|$、静态/可再生/PV/储能/VPP、外部电网，:detect\_islands），明确不以 SLACK
母线类型作代理（注释 :detect\_islands）。如实说明：\textbf{每个连通分量都生成}
\texttt{IslandInfo}（含纯 DC 岛，:detect\_islands）；无源岛（\fld{has\_generators}=false）
不删除，而是在 \texttt{AdaptiveSolver::solve} 中按死岛处理、电压置零
（\srcpath{src/power\_flow/adaptive\_solver.cpp:solve}，
\srcpath{set\_dead\_island\_zero}）。

\textbf{求解流程}（src/power\_flow/adaptive\_solver.cpp:solve）：
\begin{enumerate}
  \item 初值取母线记录（:solve），调用 \srcpath{detect\_islands}（:solve）；
  \item \textbf{单岛快路径}：不做子系统抽取，直接全系统 Newton，失败则平启动重试
        （:solve）；
  \item \textbf{多岛}：死岛（\fld{has\_generators}=false）立即置零
        （:solve，\srcpath{set\_dead\_island\_zero}）；活岛缺 AC slack 且
        \fld{enable\_auto\_swing\_selection} 时自动选摇摆母线
        （\srcpath{auto\_select\_swing\_bus}：SLACK $\to$ 最大 $P_g$ $\to$
        PV $\to$ 首母线；覆盖在抽取时提升为 SLACK 且 $\theta=0$，
        src/power\_flow/island\_detector.cpp:extract\_island\_subsystem）；
  \item 每岛 \srcpath{extract\_island\_subsystem} 重编号抽取（src/power\_flow/island\_detector.cpp:extract\_island\_subsystem，
        含开关/变压器/换流器与全部组件表），独立 Newton $+$ 平启动重试
        （src/power\_flow/adaptive\_solver.cpp:solve）；
  \item \textbf{并行}：可解岛 $>1$ 时经 \srcpath{util::ThreadPool::global()} 并发，
        每岛强制 \fld{ac\_eval\_threads}=1 防嵌套并行死锁（:solve）；单可解岛保留
        雅可比级并行（:solve）；
  \item 结果按排序母线序映射回全局向量（:map\_island\_result\_back）；
        \fld{converged} 取各岛合取、\fld{iterations} 求和、\fld{residual} 取 max
        （:solve）。
\end{enumerate}

\subsubsection*{求解参数}
\begin{paramtable}
\fld{enable\_auto\_swing\_selection} & --- & 布尔 & --- & true &
  无 slack 活岛自动选摇摆母线（src/power\_flow/adaptive\_solver.cpp:solve）\\
\fld{ac\_eval\_threads} & $n_{thr}$ & 线程 & 0（自动） & 0 &
  单可解岛的雅可比并行线程数；多岛时内部强制为 1（:solve）\\
\fld{tol}、\fld{max\_iter} & --- & --- & --- & $10^{-8}$/50 &
  转发各岛 Newton（见第~\ref{sec:newton}~章）\\
\fld{loss\_model} & --- & 枚举 & --- & Linear &
  各岛 \srcpath{make\_solver\_data} 装配损耗模型（:merge\_island\_diagnostics, :merge\_island\_diagnostics）\\
\end{paramtable}

\subsubsection*{验证与校核}
\srcpath{tests/test\_advanced\_pf.cpp}（标签 \texttt{[island]}/\texttt{[adaptive]}）：
\begin{itemize}
  \item \textbf{分岛检测}：9 母线网络断 1/2 条支路得 2/3 岛（:build\_9bus\_island\_network）；
        非连续编号回归（12 母线三馈线，断言恰为 3 岛且自适应求解后全部
        $V>0.8$，:TEST\_CASE("Island detection: non-contiguous bus IDs (case\_2\_1\_backup) → 3 islands")）；case14/case118 单岛（:TEST\_CASE）；
  \item \textbf{孤岛求解}：2 岛下 DG 岛存活、死岛电压近零、活岛
        $V\in(0.7,1.2)$（:TEST\_CASE）；混合 AC/DC 微网"并网--孤岛--重并网"
        全生命周期（:TEST\_CASE）；网状微网选择性孤岛 4 岛与互济合并
        （:TEST\_CASE）；
  \item \textbf{自适应求解}：14 母线 AC/DC 系统死母线 $|V_m|\le10^{-12}$
        （:TEST\_CASE）；case9 收敛且残差 $<10^{-4}$（:TEST\_CASE）；
        两岛集成用例功率平衡（:TEST\_CASE）；另有两节点冒烟
        （\srcpath{tests/test\_hacdcpf.cpp:TEST\_CASE}）。
\end{itemize}
